\section{Unified Mobile Private Calendar}
\label{sec:pilot-unified-mobile-calendar}

Pilot's unified mobile application now exposes the existing governed calendar authority through
the same possession-bound phone identity used by Personal Pilot.  This is not a second calendar
database and it is not a calendar running in the television.  The mother brain remains the owner
of normalized events, proposals, approvals, provider bindings and receipts; the phone is a private
review and authorization surface.åç

\subsection{User Capability}

The Personal mode contains a Calendar shortcut.  After opening it, the authenticated person can:
\begin{itemize}
  \item inspect only their own normalized calendar events, pending proposals and receipts;
  \item query free slots over a selected interval without disclosing the titles or contents of
        events occupying other times;
  \item prepare event creation, rescheduling and cancellation;
  \item specify the exact calendar, title, start, end, time zone, location, attendees, travel
        allowance and reminder interval;
  \item see whether attendee notifications are proposed before any message is sent;
  \item review conflicts and explicitly accept a conflict without revealing the conflicting
        event title;
  \item approve or reject the immutable proposal from the signed phone;
  \item execute an approved proposal only as a separate deliberate action; and
  \item inspect a normalized verified provider receipt after execution.
\end{itemize}

When no authorized calendar account is connected, Pilot remains useful for local preparation,
conflict checking and review but stops honestly before external execution.  The interface states
that the adapter is unavailable rather than representing a prepared proposal as a real event.

\subsection{Authority and Request Flow}

Each calendar operation is signed by the phone's non-exportable Ed25519 possession key and is
validated against its current mother-issued certificate and short-lived capability lease.  The
calendar scopes are separated as follows:
\begin{itemize}
  \item \texttt{calendar.read} for private snapshots and availability;
  \item \texttt{calendar.propose} for preparing an exact change;
  \item \texttt{calendar.approve} for the owner's immutable-scope decision; and
  \item \texttt{calendar.execute} for the separately authorized provider operation.
\end{itemize}

The request persona must equal the persona cryptographically bound to the phone certificate.
Consequently, a shared television request cannot silently select a person, account or calendar,
and one adult's phone cannot read or alter another adult's schedule.  Existing phones whose lease
predates these scopes must reconnect once; lease renewal cannot silently add authority.

The execution sequence is:
\begin{center}
Phone request $\rightarrow$ signed lease validation $\rightarrow$ persona calendar authority
$\rightarrow$ exact proposal hash $\rightarrow$ private decision $\rightarrow$ explicit execution
$\rightarrow$ provider verification $\rightarrow$ normalized receipt.
\end{center}

\subsection{Conflict Privacy and Retry Safety}

Conflict detection includes the selected travel allowance before the proposed start.  A conflict
record contains only an internal event reference and the occupied start and end times; it does not
contain the private title.  Pilot requires an additional explicit conflict acknowledgement before
an overlapping proposal can enter the approved state.

Every mobile preparation carries an idempotency key bound to a canonical hash of the persona,
action and complete proposed scope.  Repeating an identical request returns the original proposal.
Reusing the key after changing any field fails closed.  Executing an already completed proposal
returns the existing receipt rather than creating a second external event.

\subsection{Provider Boundary}

Persona-specific Google OAuth bindings remain on the mother brain.  Raw access and refresh tokens,
vault references and account secrets are never returned to the phone or projected onto the shared
television.  The Google Calendar creation adapter accepts both the canonical calendar proposal
shape and the earlier service-proposal shape, uses a deterministic provider event identifier for
reconciliation, and returns the provider event identifier and provider status required by the
calendar receipt authority.  Reschedule and cancellation require the exact provider event and
calendar identifiers.

Google Calendar is the first provider adapter, not the calendar architecture.  The normalized
proposal, approval and receipt contracts remain provider-independent so another authorized
calendar service can be registered without changing the Personal Pilot interface or its identity
boundary.

\subsection{Failure Semantics}

The implementation distinguishes preparation, approval and execution.  Network failure or an
unverifiable provider response moves an attempted operation to reconciliation-required state; it
does not retry a consequential provider write blindly.  Missing adapters, expired leases, revoked
phones, altered hashes, invalid attendees, invalid time zones and cross-person requests all fail
closed.  A provider effect is recorded only after an adapter returns a verified provider event
identifier.

\subsection{Verification}

Automated evidence covers private-person filtering, phone signatures and calendar scopes,
availability without title disclosure, exact attendees and calendar selection, travel-aware
conflicts, explicit conflict acceptance, idempotent preparation, changed-scope rejection,
adapter absence, verified execution receipts and duplicate-execution recovery.  At completion,
all 457 AION Fabric tests passed and the optimized Next.js application compiled and statically
generated the unified Pilot mobile route.  This is coded closure of \texttt{MOB-0504}; a real
Google account still requires the owner's normal OAuth authorization before an external event can
be created.
